%\documentclass[11pt,a4paper]{jsarticle}
\documentclass[a4paper,dvipdfmx]{jsarticle}
%
\usepackage{amsmath,amssymb}
\usepackage{bm}
\usepackage{graphicx}
\usepackage{ascmac}
\usepackage[lined,linesnumbered,boxed]{algorithm2e}
\usepackage{here}
\usepackage[subrefformat=parens]{subcaption}
\usepackage{diagbox}
\usepackage{comment} 
\SetAlgorithmName{アルゴリズム}{アルゴリズム}{アルゴリズムのリスト}
\SetKwInput{KwIn}{入力}
\SetKwInput{KwOut}{出力}
%
\setlength{\textwidth}{\fullwidth}
\setlength{\textheight}{40\baselineskip}
\addtolength{\textheight}{\topskip}
%\setlength{\voffset}{-0.2in}
%\setlength{\topmargin}{0pt}
%\setlength{\headheight}{0pt}
%\setlength{\headsep}{0pt}
\DeclareMathOperator{\err}{err}
\DeclareMathOperator{\Int}{Int}
\DeclareMathOperator{\lca}{lca}
\DeclareMathOperator{\argmin}{argmin}
\DeclareMathOperator{\BUTT}{BUTT}
\DeclareMathOperator{\Partitioning}{Partitioning}
\DeclareMathOperator{\iftext}{if}
\DeclareMathOperator{\main}{main}
\DeclareMathOperator{\IR}{IR}
\DeclareMathOperator{\SIR}{SIR}
\DeclareMathOperator{\LCAM}{LCA}
\DeclareMathOperator{\CLCAM}{ConfigLCA}
\DeclareMathOperator{\CkSS}{Config{\it k}SS}
\DeclareMathOperator{\MV}{minvalue}
\DeclareMathOperator{\kSS}{{\it k}SS}
\DeclareMathOperator{\Init}{Initialize}
\DeclareMathOperator{\key}{key}
\DeclareMathOperator{\Insert}{insert}
\DeclareMathOperator{\deletemin}{deletemin}
\DeclareMathOperator{\increasekey}{increasekey}
\DeclareMathOperator{\findmin}{findmin}
\DeclareMathOperator{\Sort}{Sort}
\DeclareMathOperator{\add}{add}
\DeclareMathOperator{\delete}{del}
\DeclareMathOperator{\NULL}{NULL}
\DeclareMathOperator{\Incresekey}{Increase{\it key}}
\DeclareMathOperator{\Changekey}{Change{\it key}}

\SetKwRepeat{Do}{do}{while}
\SetKwProg{Def}{def}{:}{end}
%
\newtheorem{theorem}{定理}[section]
\newtheorem{lemma}[theorem]{補題}
\newtheorem{corollary}[theorem]{系}
\newtheorem{proposition}[theorem]{命題}
\newtheorem{property}[theorem]{性質}
\newtheorem{remark}[theorem]{注意}
\newtheorem{example}[theorem]{例}
\newtheorem{conjecture}[theorem]{予想}
\newtheorem{exercise}[theorem]{例題}
\newtheorem{observation}[theorem]{観察}
\newtheorem{definition}[theorem]{定義}
\def\proof{\par\noindent{(証明)} \ignorespaces}
%
\title{報告書}
\author{石田　京太郎}
\date{\today}
\begin{document}
\maketitle
%
%
%\vspace{\baselineskip}
%\vspace{\baselineskip}

\section{MINCC}
グラフ$G = (X, F)$と各点対の重み$w : X \times X \rightarrow \mathbb R$が与えられたとき, 
\begin{equation}\label{eq:q_ultram}
\left|\begin{array}{rl}
\min & \displaystyle\sum_{\{i, j\} \in F \bigtriangleup \bar F} w(i,j)\\
\mbox{s.t.} & \bar G = (X, \bar F)はクリークの直和. \\
\end{array}
\right.
\end{equation}
を{\bf MINCC}と呼ぶ. ここで$\bar G = (X, \bar F)$は$G = (X, F)$を編集したグラフである. そして, 0-1行列$M  : X \times X \rightarrow \{0, 1\}$を
\begin{align}
M(i,j)=
\begin{cases}
0 &\text{if $\{i, j\} \in F$,}\\
1 &\text{otherwise}
\end{cases}
\end{align}
と定義すると, MINCCは
\begin{equation}\label{eq:q_ultram2}
\left|\begin{array}{rll}
\min & \displaystyle\sum_{i, j \in X} w(i,j)|D_{(T, l)}(i,j)- M(i,j)| & \\
\mbox{s.t.} & (T = (V, E), l)はXを葉集合とする超距離木. & \\
& D_{(T, l)}(i,j) \in \{0,1\} &(i, j \in X) \\
\end{array}
\right.
\end{equation}
と変形することができる. ここで, $D_{(T, l)}(i, j)$は$(T = (V, E), l)$における$i$-$j$間の距離である. 
\par 根付き$T = (V, E)$の点への重み付け$h: V \rightarrow \mathbb R$は, 
\begin{enumerate}
  \item $h(v) = 0$ \quad $(v \in X)$, 
  \item $h(j) \leq h(i)$ \quad $((i, j) \in E)$
\end{enumerate}
を満たすとき, {\bf 単調}と呼ばれる. 任意の単調な点重み付き根付き木$(T,h)$に対して, 
\begin{align}
  D_{(T,h)}(i, j) = h({\rm lca}_T(i, j))\quad (i, j \in X) \label{ten}
\end{align}
と定義する. ここで, ${\rm lca}_T(i, j)$は$T$中の$i$と$j$の最小共通祖先である. 式(\ref{ten})より, 問題(\ref{eq:q_ultram2})は以下の問題と等価である. 
\begin{equation}\label{eq:q_ultram3}
\left|\begin{array}{rll}
\min & \displaystyle\sum_{i, j \in V} w(i,j)|D_{(T, h)}(i,j)- M(i,j)| & \\
\mbox{s.t.} &(T = (V, E),h)は葉集合がXであるような単調な点重み付き根付き木. & \\
& h(v) \in \{0,1\} &(v \in V) \\
\end{array}
\right.
\end{equation}
問題(\ref{eq:q_ultram3})の目的関数は以下のように書き換えることができる. 
\begin{align*}
  \sum_{i, j \in V} w(i,j)|D_{(T, h)}(i,j)- M(i,j)| &=  \sum_{i, j \in V} w(i,j)|h({\rm lca}_T(i, j)) - M(i,j)| \\
&=  \sum_{v \in V} \sum_{\substack{i,j \in V \\ v = {\rm lca}_T(i, j)}} w(i,j)|h(v) - M(i,j)| \\
&=  \sum_{v \in V} \sum_{(i,j) \in d_T(v)} w(i,j)|h(v) - M(i,j)| %\\
%&=  \sum_{v \in V} \sum_{\substack{(i,j) \in d_T(v) \\ h(v) \neq M(i,j)}} w(i,j) . 
\end{align*}
ここで, $d_T(v) = \{(i, j) \mid v=\text{lca}_T(i, j), i, j \in V\}$である. 問題(\ref{eq:q_ultram3})は以下のように変形することができる. 
\begin{equation}\label{eq:q_ultram4}
\left|\begin{array}{lll}
\min & \displaystyle \sum_{v \in V} \sum_{(i,j) \in d_T(v)} w(i,j)|h(v) - M(i,j)|& \\
\mbox{s.t.} &(T = (V, E), h)は葉集合がXであるような単調な点重み付き根付き木. & \\
& h(v) \in \{0,1\} &(v \in V) \\
\end{array}
\right.
\end{equation}

\section{$\{0, 1\}$値$\ell_1$-単調回帰問題}
局所探索アルゴリズムの各反復では, $T$を固定した問題(\ref{eq:q_ultram4})を解く必要がある. この問題は根付き木上の{\bf$\{0, 1\}$値$\ell_1$-単調回帰問題}として以下のように定式化することができる.  
\begin{equation}\label{eq:q_ultram5}
\left|\begin{array}{lll}
\min & \displaystyle \sum_{v \in V} \sum_{(i,j) \in d_T(v)} w(i,j)|h(v) - M(i,j)|& \\
\mbox{s.t.} &hはTに対する単調な点重み付け. & \\
& h(v) \in \{0,1\} &(v \in V) \\
\end{array}
\right.
\end{equation}

%%%%%%%%%%%%%%%%%%%%%%%%%%%%%%%%
\begin{comment} %分かりづらくなる
$h(v) \in \{0,1\}$ $(v \in V)$, $M(i,j) \in \{0, 1\}$ $(i, j \in V)$より, 問題(\ref{eq:q_ultram5})の目的関数は, 
\begin{align*}
\sum_{v \in V} \sum_{(i,j) \in d_T(v)} w(i,j)|h(v) - M(i,j)| &=  \sum_{v \in V} \sum_{\substack{(i,j) \in d_T(v) \\ h(v) \neq M(i,j)}} w(i,j) . 
\end{align*}
と変形することができる. 
よって, 問題(\ref{eq:q_ultram5})は以下のように書き換えることができる. 
\begin{equation}\label{eq:q_ultram6}
\left|\begin{array}{lll}
\min & \displaystyle \sum_{v \in V} \sum_{\substack{(i,j) \in d_T(v) \\ h(v) \neq M(i,j)}} w(i,j)& \\
\mbox{s.t.} &hはTに対する単調な点重み付け. & \\
& h(v) \in \{0,1\} &(v \in V) \\
\end{array}
\right.
\end{equation}
\end{comment}
%%%%%%%%%%%%%%%%%%%%%%%%%%%%%%%%%

アルゴリズム~\ref{L1IR}は根付き木上の$\{0, 1\}$値$l_1$単調回帰問題に対するアルゴリズムである. ここで, アルゴリズム~\ref{L1IR}に出てくる$C(v)$, $p(v)$は, $v \in V$に対して
\begin{align*}
C(v) &= \{u : u \in V, uはvの子供\},  \\
p(v) &はvの親
\end{align*}
である. また, アルゴリズム\ref{L1IR}の2-3行目について, 
\begin{align}
Z(v,0) &= \err(v,0) + \sum_{u \in C(v)} Z(u,0),  \label{z0} \\
Z(v,1) &= \err(v,1) + \sum_{u \in C(v)}\min\{Z(u,0),Z(u,1)\} \label{z1}
\end{align}
は再帰的に定義される. $Z(v,0)$は$h(v) = 0$としたときの$v$を根とする部分木に対する部分問題の最適値であり, $Z(v,1)$は$h(v) = 1$としたときの$v$を根とする部分木に対する部分問題の最適値である. ここで, $\err$は
\begin{align}
\err (v, 0) &= \sum_{(i,j) \in d_T(v)} w(i,j)|0-M(i,j)| \notag \\
&= \sum_{\substack{(i,j) \in d_T(v) \\ M(i,j) = 1}} w(i,j), \\
\err (v, 1)  &= \sum_{(i,j) \in d_T(v)} w(i,j)|1-M(i,j)|\notag \\
&= \sum_{\substack{(i,j) \in d_T(v) \\ M(i,j) = 0}} w(i,j) 
\end{align}
と定義される. 
%%%%%%%%%%%%%%%{0,1}値l1単調回帰問題に対するアルゴリズム%%%%%%%%%%%%%%%%%%%
\DecMargin{0.5em}
\IncMargin{1.0em}
\begin{algorithm}[h]
\caption{根付き木上の${\{0, 1\}}$値$l_1$単調回帰問題に対するアルゴリズム. }
\label{L1IR}
\KwIn {根付き木$T = (V, E)$, 0-1行列$M$, 辺重み$w$, $d_T$.} 
\KwOut {$T$上の${\{0, 1\}}$値$l_1$単調回帰問題の極小最適解$h$.}
	\For{$v \in V$を$T$において後行順で}{
		$Z(v,0) \leftarrow \err(v,0) + \sum_{u \in C(v)} Z(u, 0)$\;
    		$Z(v,1) \leftarrow \err(v,1) + \sum_{u \in C(v)}\min\{Z(u, 0),Z(u, 1)\}$\;
   	 }
   	\For{$v \in V$を$T$において先行順で}{
        		\lIf{$Z(v,0) \leq Z(v, 1)$または$h(p(v)) = 0$}{
			$h(v) \leftarrow 0$
		}\lElse{
			$h(v) \leftarrow 1$
		}
	}
\end{algorithm}

\begin{example}アルゴリズム\ref{L1IR}の動作の例を示す. 入力の根付き木を図\ref{tree}とし, 0-1行列$M$を
\begin{equation}
	M =
	\begin{pmatrix}
	0 & 0 & 0 & 1 & 1\\ 
	0 & 0 & 1 & 0 & 1\\ 
	0 & 1 & 0 & 0 & 1\\ 
	1 & 0 & 0 & 0 & 1\\ 
	1 & 1 & 1 & 1 & 0\\ 
	\end{pmatrix}
, 
\end{equation}
辺重み$w$を
\begin{equation}
	w =
	\begin{pmatrix}
	0 & 2 & 1 & 3 & 4\\ 
	2 & 0 & 5 & 6 & 1\\ 
	1 & 5 & 0 & 3 & 2\\ 
	3 & 6 & 3 & 0 & 4\\ 
	4 & 1 & 2 & 4 & 0\\ 
	\end{pmatrix}
\end{equation}
とする. そして, 入力の根付き木から得られる$d_T$は以下のとおりである. 
\par [$v = 1, 2, 3, 4, 5$のとき] $d_T(v) = \{(v,v)\} $. 
\par [$v = 6$のとき] $ d_T(6) = \{(4,5)\} $. 
\par [$v = 7$のとき] $ d_T(7) = \{(1,3)\} $.  
\par [$v = 8$のとき] $ d_T(8) = \{(2,4),(2,5)\} $.  
\par [$v = 9$のとき] $ d_T(9) = \{(1,2),(1,4),(1,5),(2,3),(3,4),(3,5)\} $.  
\begin{figure}[h]
\center
\includegraphics[width=7cm]{figs/example09171.png}
\caption{入力の根付き木$T = (V, E)$.}
\label{tree}
\end{figure}
\par アルゴリズム\ref{L1IR}の2-3行目, $\err(v,0),\err(v,1)$の計算について, 例として$\err(9,0)$と$\err(9,1)$の計算される様子を以下に示し, 全ての点に対する$\err$の計算結果を図\ref{tree2}に示す. 
\begin{align*}
\err (9, 0) &= \sum_{v \in V} \sum_{\substack{(i,j) \in d_T(v) \\ h(v) \neq 0}} w(i,j) \\
&= w(1,4) + w(1,5) + w(2,3) + w(3,5) \\
&= 14,  \\
\err (9, 1) &= \sum_{v \in V} \sum_{\substack{(i,j) \in d_T(v) \\ h(v) \neq 1}} w(i,j) \\
&= w(1,2) + w(3,4)\\
&= 5. 
\end{align*}
\begin{figure}[h]
\center
\includegraphics[width=7cm]{figs/example09172.png}
\caption{各点$v$に置ける$\err$の計算結果$\{\err(v,0), \err(v,1)\}$. }
\label{tree2}
\end{figure}
\par 次にアルゴリズム\ref{L1IR}の2-3行目で$Z$が計算される様子を以下に示す. 
\par [$v = 1,2,3,4,5$のとき] $Z(v, 0)$と$Z(v, 1)$は
\begin{align*}
Z(v, 0) &= \err(v,0) + \sum_{u \in C(v)} Z(u,0) = 0, \\
Z(v, 1) &= \err(v,1) + \sum_{u \in C(v)}\min\{Z(u,0),Z(u,1)\} = 0
\end{align*}
と計算される. 
\par [$v = 6$のとき] $Z(6, 0)$と$Z(6, 1)$は
\begin{align*}
Z(6, 0) &= \err(6,0) + \sum_{u \in C(v)} Z(u,0) \\
&= 4 + 0 + 0 = 4, \\
Z(6, 1) &= \err(6,1) + \sum_{u \in C(v)}\min\{Z(u,0),Z(u,1)\} \\
&= 0 + 0 + 0 = 0. \\
\end{align*}
と計算される. 
\par [$v = 8$のとき] $Z(6, 0)$と$Z(6, 1)$は
\begin{align*}
Z(8, 0) &= \err(8,0) + \sum_{u \in C(v)} Z(u,0) \\
&= 1 + 4 + 0 = 5, \\
Z(8, 1) &= \err(8,1) + \sum_{u \in C(v)}\min\{Z(u,0),Z(u,1)\} \\
&= 6 + 0 + 0 = 6. \\
\end{align*}
と計算される. 
以下同様に$v = 7,  9$について計算した結果を図\ref{tree3}に示す. 
\begin{figure}[h]
\center
\includegraphics[width=7cm]{figs/example09173.png}
\caption{各点$v$に置ける$Z$の計算結果$\{Z(v,0), Z(v,1)\}$. }
\label{tree3}
\end{figure}
\par アルゴリズム\ref{L1IR}の5-8行目において, $T$上の${\{0, 1\}}$値$l_1$単調回帰問題の極小最適解$h$は, 先行順に
\begin{align*}
Z(v,0) \leq Z(v,1) &\Rightarrow h(v) = 0, \\
Z(v,0) > Z(v,1) &\Rightarrow h(v) = 1
\end{align*}
と決定される. しかし, 点6において$Z(v, 0) > Z(v,1)$であるが, 親である点8に対して$h(8) = 0$であるため, $h(6) = 0$となる. 最終的に出力された$h$を以下の表に示す. 
\begin{table}[H]
\caption{ $T$上の${\{0, 1\}}$値$l_1$単調回帰問題の極小最適解$h$. }
\label{ex-multi-tildez}
\begin{tabular}{l|lllllllll}
%\cline{2-10}
                           $v$  & $1$ & $2$ & $3$ & $4$ & $5$ & $6$ & $7$ & $8$ & $9$ \\ \hline
\multicolumn{1}{l|}{$h(v)$} & $0$ & $0$  & $0$  & $0$  & $0$  & $0$  & $0$  & $0$ & $1$ \\ %\hline
\end{tabular}
\centering
\end{table}
\end{example}

%\begin{comment}
\section{局所探索アルゴリズム}
アルゴリズム\ref{kSS}は, 問題(\ref{eq:q_ultram4})に対する局所探索アルゴリズムである. ここで関数$f(T,h)$を解が$(T,h)$のときの問題(\ref{eq:q_ultram4})の目的関数値とし, ${\cal N}_k(T)$を$T$の$k$SS近傍とする. アルゴリズム\ref{L1IR}は$d_T$と$\err$の計算をあらかじめ行っておくことでO($n$)で実行可能である. 次に6行目の$d_{T'}$と$\err$の計算については, 水越修論{\bf 補題} 3.13 総和行列の計算時間の証明の部分より, O($k$)で実行できそうであることが分かった. よって, 局所探索アルゴリズム1反復当たりの計算時間はO$(n \times k{\rm SS近傍の大きさ})$になりそうである. 
\begin{algorithm}[h]
\caption{局所探索アルゴリズム. }
\label{kSS}
\KwIn {根付き木$T = (V,E)$, 0-1行列$M$, 辺重み$w$. }
\KwOut {$l_1$最良近似超距離木問題の局所最適解となる根付き木$T = (V, E)$.}
$d_T$と$\err$の計算\;
$T$上の$\{0,1\}$値$\ell_1$単調回帰問題の極小最適解$h$を求める\;
$\MV \leftarrow f(T, h)$\;
\Repeat{$\MV$が更新されなくなる}{
	\For{$T' \in {\cal N}_k(T)$}{
		$d_{T'}$と$\err$の計算\;
		$T'$上の$\{0,1\}$値$\ell_1$単調回帰問題の極小最適解$h'$を求める\;
		\If{$f(T', h') < \MV$}{
			$\MV \leftarrow f(T', h')$\;
			$T \leftarrow T'$\;
		}
	}
}
\end{algorithm}
%%%%%%%%%%%%%%%%%%%%%%%%%%%%%%%%%%%%%%%%%%%%%%%%%%%%%%%%%%%%%%%%%%%%%%%%%%%%%%%%%%%%%%%%%%%%%%%%%%%%%%%%%%%%%%%%%%%%%%%%%%%%%%%%%%%%%%%%%%%%%%%%%%%%%%%%%%%%%%%%%%%%%%%%%%%%%%%%%%%%%%%%%%%
%%%%%%%%%%%%%%%%%%%%%%%%%%%%%%%%%%%%%%%%%%%%%%%%%%%%%%%%%%%%%%%%%%%%%%%%%%%%%%%%%%%%%%%%%%%%%%%%%%%%%%%%%%%%%%%%%%%%%%%%%%%%%%%%%%%%%%%%%%%%%%%%%%%%%%%%%%%%%%%%%%%%%%%%%%%%%%%%%%%%%%%%%%%%
\section{次回の目標}
\begin{itemize}
	\item 局所探索アルゴリズムの開発. 
\end{itemize}
%\end{comment}

\end{document}
